\section{Los tres exponentes en un soporte}
\begin{proposicion}\label{prop:s}
Bajo los márgenes de la sección 1, para \(5\le v\le8\log n\), sea \(c_s\) el número
de componentes del grafo bipartito de los sitios de \(U\) y sus fibras de las cuatro
familias. Sea \(\beta_U\) su número total de ciclos independientes, y
\(d_i=\max_{f\ni i}|f\cap U|\). Para el aproximante de la proposición~\ref{prop:m},
\[
 \left|\Delta_U(\widehat F_{U,q}-\Lambda)\right|
 \le C^v(Cv)^{c_s+1}j^{-p}\prod_{i\in U}d_i,\qquad
 p=\begin{cases}
 v+c_s+1 & \beta_U=0,\\
 v+c_s & \beta_U=1,\\
 v+c_s-1 & \beta_U\ge2 .
 \end{cases}
\]
La cota se aplica a los soportes pertinentes, en los que cada sitio participa en un
conflicto activable. Filas y columnas permanecen en el grafo aunque se apague una diagonal.
\end{proposicion}
\begin{proof}
Sea \(s=1/j\) en el punto físico y libérense los campos \(\zeta_i\) de cada sitio
después de extraer el mismo campo afín en todas las caras. El balanceo de \pruebaen{B01} y
\pruebaen{B02} admite las deformaciones de medida de tamaño \(O(v|s|)\).
Por \(\sum_{i\in f}1/d_i\le1\), la silla local y su determinante en \pruebaen{G02}
son holomorfos en \(|s|<\kappa/v,\ |\zeta_i|<\kappa'/d_i\). Las diferencias mixtas
de los pesos en \pruebaen{G03} y la inducción ocupado/vacío en \pruebaen{G04} dan
un logaritmo del conteo aproximado \(O(v)\) en esa caja. Para la inducción basta
\(\prod_D(1+|e^{a_D}-1|)\le\exp(\sum_D|a_D|)\), por la serie exponencial.
La rama del predictor coincide con la solución pequeña de \pruebaen{N02}.
El cofactor común de \pruebaen{U07} hace coincidir las caras en \(\zeta_i=0\).
El alternado \(H(s,\zeta)\) está pues acotado por \(C^v\) y es divisible por todos los
\(\zeta_i\).

Para contar potencias de \(s\), se toma un grupo por componente. Un monomio conexo de coste
\(w\), con interacciones de costes \(\ell_a\) que tocan \(h_a\) grupos distintos, ocupa
\(w+1-\sum_a(\ell_a+1-h_a)-L_I\) grupos, donde \(L_I\ge0\) es el número de ciclos del
grafo de incidencias del monomio. Los déficits son no negativos por las dos fórmulas del
exterior de \pruebaen{B01} y su desarrollo en \pruebaen{P03}. Los coeficientes sobre grupos
separados se factorizan antes de tomar logaritmos. Al alternar, un grupo ausente da cero.
Por ello desaparecen todos los costes inferiores a \(c_s-1\).

En el bosque la expansión conexa centrada de \pruebaen{P04} cancela las contracciones
gaussianas del único ciclo ordinario contra el medio logaritmo del determinante.
La única interacción excepcional que aún podría alcanzar \(w\) grupos contiene dos
inserciones con la misma amplitud. Podar una rama sustituye sus hojas por sus medias y
una derivada en el campo del grupo de enlace en ambas expresiones. En la raíz, con las
ocupaciones binarias \(Y_1,Y_2\), \pruebaen{X02} obtiene que el residuo tras la traza es
\(4(1-\mathbb E Y_1)(1-\mathbb E Y_2)\operatorname{Cov}(Y_1,Y_2)\).
Para alcanzar el soporte máximo esas fibras están en grupos distintos, independientes
también al variar sus campos. El residuo y sus derivadas se anulan. Quedan a lo sumo
\(w-1\) grupos, de modo que el primer coste es al menos \(c_s+1\).

Si hay un solo ciclo, para alcanzar el coste \(c_s-1\) todas las ramas fuera de su
componente se podan del mismo modo. Si \(k_i\) es el número de factores incidentes,
la razón ocupado/vacío en el producto de creencias es
\(x_i^{k_i}/[\prod_{f\ni i}q_f\,(x_i/(1-x_i))^{k_i-1}]=\zeta_i\) por estación.
El peso basal es \(Z_{\rm Bethe}^{-1}\). Las transiciones binarias del ciclo dan
\(Z/Z_{\rm Bethe}=1+\theta\). En el complemento del modo nulo axial, el lema
del determinante de \pruebaen{U07} reduce el cofactor a
\(\det[(-H_B''-R'')V_0]\prod_f\bigl(q_f/\prod_{i\in f}(1-x_i)\bigr)\), con el producto
de cocientes tomado sobre todas las fibras y \(V_0=\diag(x_i(1-x_i))\).
A reacción exterior cero, la identidad de incidencias de \pruebaen{Z01} da
\((1-\theta)^2\), un factor por orientación. El error es \(\log(1-\theta^2)\),
con \(\theta=(-1)^m\prod_{i\in\mathcal C}x_i/(1-x_i)\) para sus \(m\) sitios.
Cada \(x_i\) es divisible por \(\zeta_i\). Las derivadas de campo
\(\zeta_i\partial_{\zeta_i}\) conservan la doble multiplicidad. \pruebaen{Z02} y
\pruebaen{P05} dan así un campo adicional de cada sitio del ciclo en el término crítico.
Cuando hay dos o más ciclos basta la cuenta general de conexidad.

El paso al alcance creciente es \pruebaen{X03} y \pruebaen{V07}. Cambiar el orden alto
por \(q=4v+12\) conserva todos los grados leídos, ya que el cambio de germen está en
\((s^{q-1})\). Por identidad holomorfa se prolongan a toda la caja las anulaciones y
las dobles divisibilidades probadas inicialmente para cada soporte. Cauchy en \(s\)
desde el primer coste \(w\) da \(C^v(Cv/j)^w\). Schwarz sucesivo en activaciones
aporta \(\prod_i(Cd_i/j)\). En el coste crítico unicíclico los campos adicionales
cuestan a lo sumo \(\sum_i Cd_i/j\le Cv^2/j\) en el punto físico.
Así salen los tres exponentes con la constante anunciada. El tamaño del permanente
físico sólo interviene después, en la comparación real de la proposición~\ref{prop:m}.
\end{proof}
Para \(v\le4\) se usa la versión física a grado fijo de \pruebaen{P05}, incluida su
proyección del resto para cada tamaño entero. Los casos geométricos se especifican
en B.1 y en la pieza V05 del archivo de fuentes.
